\chapter{Application Requirements and Specification}
\label{chap:requirements}

\par This chapter defines the application that constitutes the practical contribution of the thesis: a web demonstration in which the user submits captured network traffic and receives per-flow intrusion-detection verdicts from the model of Chapter~\ref{chap:ml_methodology}.

\section{Project Goal and Scope}
\label{sec:project_goal}

\par The application is a real-time intrusion detection demonstration for IoT networks, built around the model of Chapter~\ref{chap:ml_methodology}. It does not train models: it consumes the artefacts produced by the offline training pipeline (the fitted scaler, the label encoder, and the trained weights) and serves them. Concretely, the application loads those artefacts, accepts either a pre-extracted feature CSV in CIC format or a raw packet capture (from which the project's own flow extractor derives the features), runs the trained model, and presents per-flow verdicts and an aggregated summary in the browser.

\par The system targets an interactive, single-request setting: the user supplies a sample of traffic and reviews the verdict on screen. Line-rate gateway protection (sub-millisecond per-packet inference at multi-gigabit throughput) is out of scope; this narrower sense of \enquote{real time} sets the latency thresholds below. Both classification granularities are exposed: the 2-class gate an inline deployment would use, and the 8-class attack-family view an analyst would consume.

\section{Functional Requirements}
\label{sec:functional_requirements}

\par The core requirement (FR-D1) is per-flow classification of submitted traffic: the demo accepts either a CSV in the CIC feature format or a raw packet capture, and returns a per-flow verdict table with an aggregated summary. For packet captures the features are derived by the project's own extractor, which reproduces the CIC-IoT-2023 dataset's packet-windowing method so that live features match the training distribution (Section~\ref{sec:cicflowmeter_design}). The remaining requirements support this core. Model-variant selection (FR-D2) lists every (split, granularity) pair with trained artefacts on disk and lets the user switch without restarting. The output schema is deterministic (FR-D3): a summary line with the flow count and a per-label breakdown, followed by a table of row index, predicted label, and confidence. Input is self-validating (FR-D4): a CSV missing expected columns produces an error naming them, and a capture from which no flow can be extracted produces a readable error rather than crashing the request.

\par Raw traffic is never persisted (FR-D5): uploads are processed in a temporary location, removed when the request ends, and never logged; only derived result metadata (model, predicted label, confidence, mode, timestamp) may be stored for an authenticated user. The application is deployed as a permanent public service (FR-D6) at a fixed HTTPS URL with no presenter-side configuration: the backend as a Docker container on Hugging Face Spaces and the React frontend on Vercel. Because the demonstration is publicly reachable, the deployed frontend requires authentication and keeps a per-user history of past classifications (FR-D7, metadata only per FR-D5), so that each user's history stays private to them, with isolation enforced at the storage layer rather than in client code.

\section{Non-Functional Requirements}
\label{sec:non_functional_requirements}

\par End-to-end per-flow latency through the demo (CSV parsing, scaling, forward pass, result formatting) must stay below 100~ms at p95 on a consumer laptop CPU for batches of up to 1024 flows (NFR-1); packet-capture parsing is excluded because feature extraction dominates that path, and capture-to-verdict latency is reported separately. For classification performance (NFR-2), the served model must reach at least 0.85 weighted-F1 on the held-out binary test set; the harder 8-class task is a secondary objective whose weighted-F1 is expected to be lower, since the transport-layer features cannot fully separate the application-layer families.

\par The application must run on commodity hardware (NFR-3): CPU-only, without a GPU dependency, on both Windows and Linux, with the backend packageable as a Docker image. Operation must stay simple (NFR-4): a permanent HTTPS URL, a single command for local development, and no database, message queue, or dedicated model server, since the backend loads artefacts directly into the application process. Reporting must be faithful (NFR-5): per-class metrics always accompany aggregates so majority-class dominance cannot mask weak minority classes, and the random split's susceptibility to within-session redundancy is stated plainly rather than hidden behind a high aggregate figure.

\section{Use Cases}
\label{sec:use_cases}

\par The application has a single actor, the \emph{operator}: the person who opens the web interface, selects a model variant, uploads a packet capture or feature CSV, and reads the per-flow verdicts. The principal interactions are summarised in Table~\ref{tab:use_cases}; the core classification path is then narrated as the representative flow.

\begin{table}[htbp]
\begin{center}
\begin{tabular}{|p{34pt}|p{92pt}|p{95pt}|p{125pt}|}
\hline
\textbf{ID} & \textbf{Actor} & \textbf{Trigger} & \textbf{Outcome} \\
\hline
\hline UC-1 & Operator & Uploads a packet capture & Per-flow verdict table and summary; temporary capture removed afterwards. \\
\hline UC-2 & Operator & Uploads a pre-extracted CIC CSV & Per-flow verdict table; missing columns reported as a readable error. \\
\hline UC-3 & Operator & Selects a different model card & Result screen re-rendered under the new variant with no restart or reload. \\
\hline UC-4 & Operator & Opens the public NEURAL IDS URL & Demo reachable from any network at a permanent URL with no presenter-side setup. \\
\hline
\end{tabular}
\end{center}
\caption{Principal use cases with their actors, triggers, and outcomes.}
\label{tab:use_cases}
\end{table}

\par In the core flow (UC-1), the application stores the uploaded capture temporarily, runs the flow extractor to obtain the features, applies the trained scaler and model, renders the summary and per-flow verdict table, and removes the temporary files; if flow extraction fails, its error output is returned with an empty table. The CSV path (UC-2) is identical except that it validates the CIC column set instead of running flow extraction.

\section{System Specification}
\label{sec:system_specification}

\par This section fixes the high-level architecture and the technology choices; the detailed design follows in Chapter~\ref{chap:design_implementation}.

\subsection{High-Level Architecture}
\label{subsec:high_level_architecture}

\par The system consists of two pipelines that meet at a small set of serialised artefacts. The training pipeline runs offline and produces those artefacts. The serving pipeline is the live demo, which loads them and answers classification requests. Figure~\ref{fig:system_architecture} shows the data flow.

\begin{figure}[htbp]
    \centering
    \resizebox{0.88\textwidth}{!}{%
    \begin{tikzpicture}[
        node distance=0.6cm and 0.8cm,
        layer/.style={
            rectangle,
            rounded corners=4pt,
            draw=black!70,
            thick,
            minimum width=2.6cm,
            minimum height=1.0cm,
            text width=2.5cm,
            align=center,
            font=\small
        },
        train/.style={layer, fill=blue!8},
        artefact/.style={layer, fill=yellow!18, minimum width=2.4cm},
        serve/.style={layer, fill=green!10},
        arrow/.style={
            ->,
            >=latex,
            thick,
            black!60
        }
    ]
    % Training pipeline (top row)
    \node[train] (raw)        {Raw CIC\\captures};
    \node[train, right=of raw]        (parquet)   {Cleaned\\columnar\\intermediate};
    \node[train, right=of parquet]    (split)     {Split\\(random\\stratified)};
    \node[train, right=of split]      (train)     {Model\\training};

    % Artefact layer (middle)
    \node[artefact, below=of train, yshift=-0.2cm] (scaler)  {Fitted\\scaler};
    \node[artefact, left=of scaler]                (encoder) {Fitted\\encoder};
    \node[artefact, left=of encoder]               (weights) {Trained\\weights};

    % Serving pipeline (bottom row)
    \node[serve, below=of weights, yshift=-0.2cm]  (pcap)    {Capture /\\CSV upload};
    \node[serve, right=of pcap]                    (cic)     {Flow\\extractor};
    \node[serve, right=of cic]                     (predict) {Inference\\component};
    \node[serve, right=of predict]                 (webui)  {Web UI};

    % Arrows: training row
    \draw[arrow] (raw) -- (parquet);
    \draw[arrow] (parquet) -- (split);
    \draw[arrow] (split) -- (train);

    % Train -> artefacts
    \draw[arrow] (train.south) -- (scaler.north);
    \draw[arrow] (train.south) -- (encoder.north);
    \draw[arrow] (train.south) -- (weights.north);

    % Artefacts -> predictor
    \draw[arrow] (scaler.south) -- (predict.north);
    \draw[arrow] (encoder.south) -- (predict.north);
    \draw[arrow] (weights.south) -- (predict.north);

    % Serving row
    \draw[arrow] (pcap) -- (cic);
    \draw[arrow] (cic) -- (predict);
    \draw[arrow] (predict) -- (webui);

    \end{tikzpicture}%
    }
    \caption{High-level architecture of the RT-IDS application. The training pipeline (blue) produces three serialised artefacts (yellow), which the serving pipeline (green) loads to answer per-flow classification requests.}
    \label{fig:system_architecture}
\end{figure}

\subsection{Technology Stack and Rationale}
\label{subsec:tech_stack}

\par The technology choices follow from the requirements. Features are scaled with the robust median/IQR scaler and labels use a reversible integer encoder; imbalance is handled by class-weighted cross-entropy (Section~\ref{subsec:class_weighted_ce}). The served model is the MLP of Section~\ref{sec:model_architecture}, with the Random Forest baseline trained and evaluated offline only. Packet captures are converted to features by the project's own \texttt{dpkt}-based extractor, which reproduces the windowing method the CIC-IoT-2023 dataset was itself built with; CICFlowMeter is deliberately not used, because its feature definitions differ from the dataset's (Section~\ref{sec:cicflowmeter_design}). The demo pairs a FastAPI backend with the NEURAL IDS React dashboard, containerised on Hugging Face Spaces and served from Vercel respectively, which yields a permanent HTTPS URL (NFR-3, NFR-4). The backend loads the artefacts directly into the application process, so no separate model server is needed.

\subsection{Constraints and Assumptions}
\label{subsec:constraints}

\par Several constraints bound the work. The 39-feature public CSV release of CICIoT2023 is the only source of training data; the 46-feature variant used by some published comparisons was not available from the official distribution, so numerical results are not directly comparable across the two feature sets (Section~\ref{subsec:features}). The dataset comes from a smart-home testbed, so generalisation to industrial, medical, or vehicular networks is not claimed, and its header-and-metadata features mean detection of application-layer web attacks is expected to be weak, a property of the feature set discussed in the results chapter. The serving target is a single-container backend with a static frontend, so distributed serving, load balancing, and high availability are out of scope. Because the demonstration is deployed as a public service, the per-user history (FR-D7) must be kept private to each user, which requires authentication; this is provided by Supabase email/password, chosen because it bundles managed authentication with a Postgres database and row-level security, so per-user isolation is enforced at the data layer without a bespoke authentication backend. Finally, the deployed model is the one trained offline; the demo does not adapt to feedback during a session.
